% GENERATED FILE. Edit canonical lessons, then run npm run generate:books.
% canonical-source-hash: fnv1a64:c5a11b675bf07e21
% canonical-lessons: ZH-C58-shouji, ZH-W58-ji-machine, ZH-C58-bao, ZH-C58-diannao, ZH-C58-beizi, ZH-C58-yifu

\chapter{Things You Carry}
\label{ch:zh-things-you-carry}
\hlchaptermodality{58}

\begin{chapteropening}
\textbf{By the end of this chapter:} \emph{I can name my phone, my bag, my computer, a cup and my clothes.}

Complete the last lesson of chapter 58: I can name my phone, my bag, my computer, a cup and my clothes.
\end{chapteropening}

\section[\texorpdfstring{\zh{手机} (shǒujī)}{手机 (shǒujī)}]{\zh{手机} (shǒujī) — mobile phone}

\label{lesson:ZH-C58-shouji}



\begin{quote}
Before the new one: say the Mandarin for chair, then the Mandarin for here.
\end{quote}

\subsection*{You'll want to know: \zh{手机}}
\textbf{\zh{手机}} — \emph{shǒujī} — \textquotedblleft{}mobile phone\textquotedblright{}.

\begin{quote}
\textbf{\zh{这是我的手机}} \emph{zhè shì wǒ de shǒujī} — This is my phone.
\end{quote}

The new character, in pieces: \zh{机} — the tree piece on the left, the one at the top left of \zh{想}, and on the right a two-stroke piece read jǐ, as in the word for how many. The next lesson puts it in your hand.

\begin{cousinweb}
\textbf{\zh{手机}} (\emph{shǒujī}) is built from \zh{手} 'hand' and \zh{机} 'machine'. The second is inherited from Middle Chinese \emph{kj+j} (Guangyun, in Baxter's transcription), which the dictionary explains as the catch that sets off a crossbow. From that catch to any machine is the usual explanation, but the step is not given there.
\end{cousinweb}

\begin{grammarlens}[title={What you've built}]
One more word for everyday answers.
\end{grammarlens}

\subsection*{Your turn}
\begin{itemize}
\raggedright
  \item \emph{Say it:} \emph{shǒujī}
  \item \emph{Say it:} \emph{shǒujī}, once more
  \item \emph{Say it:} \emph{zhè shì wǒ de shǒujī} — This is my phone.
  \item \emph{Recall:} say the Mandarin for chair, then the Mandarin for here, then say \emph{shǒujī} again
\end{itemize}

\begin{tcolorbox}[breakable,colback=teal!4,colframe=teal!35!black,title={Before you move on}]
What does \zh{手机} mean? (\textquotedblleft{}Mobile phone\textquotedblright{}.) Say it once more.
\end{tcolorbox}
\section[\texorpdfstring{\zh{机} (jī)}{机 (jī)}]{\zh{机} — the tree piece, then jǐ}

\label{lesson:ZH-W58-ji-machine}



\begin{quote}
The word first, then the character inside it.

\begin{itemize}
\raggedright
  \item \emph{Recall:} say \emph{shǒujī} — \textbf{R1}, one lesson back, before it decays
\end{itemize}

\begin{quote}
\zh{手} \zh{想}
\end{quote}

\emph{Read it:} each one aloud, with nothing said first — \emph{shǒu}, \emph{xiǎng}
\end{quote}

\subsection*{Script — the tree piece, then jǐ}
\hlblockfigure{\detokenize{figures/ZH-W58-ji-machine-filmstrip.pdf}}{How \zh{机} is written, stroke by stroke}

\begin{quote}
\zh{机}
\end{quote}

\textbf{Two pieces, left before right.} Count the pieces, not the strokes.

\textbf{The tree piece, four strokes}, narrow on the left, as at the top left of \zh{想}: a horizontal, a vertical down through it, a left-falling stroke, and a dot.

\textbf{The right piece, two strokes, read jǐ}, as in the word for how many: a left-falling stroke down the left side, then one stroke that runs across the top, turns down, bends right along the bottom and hooks up.

4 + 2 = six.

\subsection*{Your turn}
\begin{itemize}
\raggedright
  \item \emph{Write it:} \zh{机}, all six, with the model in front of you
  \item \emph{Cover:} write it again from memory, and count the pieces, not the strokes
\end{itemize}

\begin{tcolorbox}[breakable,colback=teal!4,colframe=teal!35!black,title={Before you move on}]
How many pieces in \zh{机}? (\textbf{Two} — the tree piece, then jǐ.) Which character you know has the same tree piece? (\textbf{\zh{想}}.) How does the last stroke end? (\textbf{It runs along the bottom and hooks up}.)

Source: \href{https://raw.githubusercontent.com/chanind/hanzi-writer-data/68d10a4b21150cae5e1ebbd223eed289cf32d90c/data/\%E6\%9C\%BA.json}{Hanzi Writer Data \zh{机}, pinned stroke paths and medians}.
\end{tcolorbox}
\section[bāo]{bāo — bag}

\label{lesson:ZH-C58-bao}



\begin{quote}
Before the new one: say the Mandarin for here, then the Mandarin for mobile phone.
\end{quote}

\subsection*{You'll want to know: bāo}
\emph{bāo} — \textquotedblleft{}bag\textquotedblright{}.

\begin{quote}
\emph{wǒ xiǎng mǎi yí ge bāo} — I'd like to buy a bag.
\end{quote}

Its characters come later; for now it is a word for the ear.

\begin{cousinweb}
\emph{bāo} is inherited from Middle Chinese \emph{paew} (Guangyun, in Baxter's transcription), glossed 'to wrap'. The sense 'bag' is not given there; a thing that wraps is the likely step, but it is not settled.
\end{cousinweb}

\begin{grammarlens}[title={What you've built}]
One more word for everyday answers.
\end{grammarlens}

\subsection*{Your turn}
\begin{itemize}
\raggedright
  \item \emph{Say it:} \emph{bāo}
  \item \emph{Say it:} \emph{bāo}, once more
  \item \emph{Say it:} \emph{wǒ xiǎng mǎi yí ge bāo} — I'd like to buy a bag.
  \item \emph{Recall:} say the Mandarin for here, then the Mandarin for mobile phone, then say \emph{bāo} again
\end{itemize}

\begin{tcolorbox}[breakable,colback=teal!4,colframe=teal!35!black,title={Before you move on}]
What does bāo mean? (\textquotedblleft{}Bag\textquotedblright{}.) Say it once more.
\end{tcolorbox}
\section[diànnǎo]{diànnǎo — computer}

\label{lesson:ZH-C58-diannao}



\begin{quote}
Before the new one: say the Mandarin for mobile phone, then the Mandarin for bag.
\end{quote}

\subsection*{You'll want to know: diànnǎo}
\emph{diànnǎo} — \textquotedblleft{}computer\textquotedblright{}.

\begin{quote}
\emph{wǒ xǐhuan wán diànnǎo} — I like playing on the computer.
\end{quote}

Its characters come later; for now it is a word for the ear.

\begin{cousinweb}
\emph{diànnǎo} is built from diàn 'lightning', now 'electricity', and nǎo 'brain', both inherited from Middle Chinese, \emph{denH} and \emph{nawX} (Guangyun, in Baxter's transcription): an electric brain.
\end{cousinweb}

\begin{grammarlens}[title={What you've built}]
One more word for everyday answers.
\end{grammarlens}

\subsection*{Your turn}
\begin{itemize}
\raggedright
  \item \emph{Say it:} \emph{diànnǎo}
  \item \emph{Say it:} \emph{diànnǎo}, once more
  \item \emph{Say it:} \emph{wǒ xǐhuan wán diànnǎo} — I like playing on the computer.
  \item \emph{Recall:} say the Mandarin for mobile phone, then the Mandarin for bag, then say \emph{diànnǎo} again
\end{itemize}

\begin{tcolorbox}[breakable,colback=teal!4,colframe=teal!35!black,title={Before you move on}]
What does diànnǎo mean? (\textquotedblleft{}Computer\textquotedblright{}.) Say it once more.
\end{tcolorbox}
\section[bēizi]{bēizi — cup; glass}

\label{lesson:ZH-C58-beizi}



\begin{quote}
Before the new one: say the Mandarin for bag, then the Mandarin for computer.
\end{quote}

\subsection*{You'll want to know: bēizi}
\emph{bēizi} — \textquotedblleft{}cup; glass\textquotedblright{}.

\begin{quote}
\emph{zhè shì nǐ de bēizi ma} — Is this your cup?
\end{quote}

Its characters come later; for now it is a word for the ear.

\begin{cousinweb}
\emph{bēizi} is built from bēi and the noun ending zi, both inherited from Middle Chinese, \emph{pwoj} and \emph{tsiX} (Guangyun, in Baxter's transcription). The dictionary explains the first, through an older book, as a kind of vessel.
\end{cousinweb}

\begin{grammarlens}[title={What you've built}]
One more word for everyday answers.
\end{grammarlens}

\subsection*{Your turn}
\begin{itemize}
\raggedright
  \item \emph{Say it:} \emph{bēizi}
  \item \emph{Say it:} \emph{bēizi}, once more
  \item \emph{Say it:} \emph{zhè shì nǐ de bēizi ma} — Is this your cup?
  \item \emph{Recall:} say the Mandarin for bag, then the Mandarin for computer, then say \emph{bēizi} again
\end{itemize}

\begin{tcolorbox}[breakable,colback=teal!4,colframe=teal!35!black,title={Before you move on}]
What does bēizi mean? (\textquotedblleft{}Cup; glass\textquotedblright{}.) Say it once more.
\end{tcolorbox}
\section[yīfu]{yīfu — clothes}

\label{lesson:ZH-C58-yifu}



\begin{quote}
Before the new one: say the Mandarin for computer, then the Mandarin for cup.
\end{quote}

\subsection*{You'll want to know: yīfu}
\emph{yīfu} — \textquotedblleft{}clothes\textquotedblright{}.

\begin{quote}
\emph{wǒ xǐhuan mǎi yīfu} — I love buying clothes.
\end{quote}

Its characters come later; for now it is a word for the ear.

\begin{cousinweb}
\emph{yīfu} is built from yī and fu, both inherited from Middle Chinese, \emph{'j+j} and \emph{bjuwk} (Guangyun, in Baxter's transcription). The dictionary glosses the first as the garment worn on the upper body, and lists 'clothes' among the senses of the second.
\end{cousinweb}

\begin{grammarlens}[title={What you've built}]
One more word for everyday answers.
\end{grammarlens}

\subsection*{Your turn}
\begin{itemize}
\raggedright
  \item \emph{Say it:} \emph{yīfu}
  \item \emph{Say it:} \emph{yīfu}, once more
  \item \emph{Say it:} \emph{wǒ xǐhuan mǎi yīfu} — I love buying clothes.
  \item \emph{Recall:} say the Mandarin for computer, then the Mandarin for cup, then say \emph{yīfu} again
\end{itemize}

\begin{tcolorbox}[breakable,colback=teal!4,colframe=teal!35!black,title={Before you move on}]
What does yīfu mean? (\textquotedblleft{}Clothes\textquotedblright{}.) Say it once more.
\end{tcolorbox}
